% ============================================================
% TLC State Diagram — paste into your report
% Requires: \usepackage{tikz}
%           \usetikzlibrary{automata, positioning, arrows.meta}
% ============================================================
\begin{figure}[htbp]
    \centering
    \begin{tikzpicture}[->, >=Stealth, shorten >=1pt, auto, node distance=4.5cm,
            semithick, every state/.style={minimum size=1.8cm, font=\large}]

        % --- States ---
        \node[state, fill=green!15]   (G) {$\mathbf{G}$};
        \node[state, fill=yellow!25]  (Y) [right of=G] {$\mathbf{Y}$};
        \node[state, fill=red!15]     (R) [right of=Y] {$\mathbf{R}$};

        % --- Transitions ---
        \path (G) edge [bend left=25]  node {\texttt{request}} (Y)
        (Y) edge [bend left=25]  node {5\,s timeout}     (R);
        \draw[->] (R.south west) .. controls +(0,-1.5) and +(0,-1.5) ..
        node[below, pos=0.5] {10\,s timeout} (G.south east);

        % --- Output labels ---
        \node[below=1.8cm of G, align=center, font=\footnotesize] {
            \underline{State G}\\[2pt]
            Veh: \textcolor{green!60!black}{\textbf{GREEN}}\\[2pt]
            Ped: \textcolor{red}{\textbf{RED}}\\[2pt]
            \texttt{output = 10001}
        };
        \node[below=1.8cm of Y, align=center, font=\footnotesize] {
            \underline{State Y}\\[2pt]
            Veh: \textcolor{yellow!70!black}{\textbf{YELLOW}}\\[2pt]
            Ped: \textcolor{red}{\textbf{RED}}\\[2pt]
            \texttt{output = 01001}
        };
        \node[below=1.8cm of R, align=center, font=\footnotesize] {
            \underline{State R}\\[2pt]
            Veh: \textcolor{red}{\textbf{RED}}\\[2pt]
            Ped: \textcolor{green!60!black}{\textbf{GREEN}}\\[2pt]
            \texttt{output = 00110}
        };

    \end{tikzpicture}
    \caption{State diagram for the traffic light controller FSM. The 5-bit output vector
        is \texttt{\{veh\_G, veh\_Y, veh\_R, ped\_G, ped\_R\}}. \texttt{reset} (active LOW, asynchronous) returns to state G from any state.}
    \label{fig:state-diagram}
\end{figure}